\chapter{优化量推导与数值实现}

\section{时刻级残差函数}
在单时刻估计中，定义第 $i$ 个信标残差
\[
r_i(\bm p)=d_i(\bm p)-\hat d_i
=\sqrt{(x-u_i)^2+(y-v_i)^2}-\hat d_i.
\]
目标函数为
\[
F(\bm p)=\sum_i \omega_i r_i(\bm p)^2.
\]

\section{梯度表达}
记 $\Delta x_i=x-u_i,\ \Delta y_i=y-v_i,\ d_i=\sqrt{\Delta x_i^2+\Delta y_i^2}$，则
\[
\frac{\partial r_i}{\partial x}=\frac{\Delta x_i}{d_i},\qquad
\frac{\partial r_i}{\partial y}=\frac{\Delta y_i}{d_i}.
\]
由链式法则有
\begin{equation}
\nabla F(\bm p)=2\sum_i \omega_i r_i
\begin{bmatrix}
\Delta x_i/d_i\\
\Delta y_i/d_i
\end{bmatrix}.
\end{equation}

\section{近似海森矩阵}
在高斯牛顿近似下，忽略二阶残差项，可写
\begin{equation}
\bm H \approx 2\sum_i \omega_i \bm j_i\bm j_i^{\T},
\quad
\bm j_i=
\begin{bmatrix}
\Delta x_i/d_i\\
\Delta y_i/d_i
\end{bmatrix}.
\end{equation}
该近似在残差较小时具有良好数值稳定性。

\section{滑动窗口展开}
窗口长度为 $w$ 时，优化变量为
\[
\bm X=[\bm p_{t-w+1}^{\T},\dots,\bm p_t^{\T}]^{\T}\in\R^{2w}.
\]
对应法方程呈稀疏块带状结构。观测项贡献块对角矩阵，运动约束贡献相邻块耦合项，适合使用稀疏线性代数求解。

\section{鲁棒迭代重加权}
在第 $k$ 次迭代：
\begin{enumerate}
  \item 用当前轨迹计算残差 $r_i^{(k)}$；
  \item 根据 Huber 规则更新权重 $\omega_i^{(k)}$；
  \item 构造并求解线性化增量方程；
  \item 更新轨迹并判断收敛。
\end{enumerate}
收敛条件可采用
\[
\frac{\|\Delta \bm X\|_2}{\|\bm X\|_2}<10^{-5}.
\]

\section{实现细节}
为保证实时性，本文实现采用以下策略：
\begin{itemize}
  \item 只保留窗口内变量，避免全局重算；
  \item 对极端 RSSI 值做截断，防止数值爆炸；
  \item 对可行域投影使用最近栅格近似，降低计算量；
  \item 每帧最多迭代固定步数，保证时延上界。
\end{itemize}
